\documentclass[11pt,a4paper]{article}

\usepackage[utf8]{inputenc}
\usepackage[T1]{fontenc}
\usepackage[french]{babel}
\usepackage[a4paper,margin=2.2cm]{geometry}
\usepackage{booktabs}
\usepackage{array}
\usepackage{longtable}
\usepackage{tabularx}
\usepackage{enumitem}
\usepackage{xcolor}
\usepackage{hyperref}
\usepackage{listings}
\usepackage{float}
\usepackage{titlesec}
\usepackage{fancyhdr}
\usepackage{amssymb}

\definecolor{dressmeblue}{RGB}{32,79,130}
\definecolor{dressmegray}{RGB}{90,90,90}
\definecolor{lightgray}{RGB}{245,245,245}
\definecolor{lightblue}{RGB}{235,243,252}
\definecolor{warning}{RGB}{150,80,0}
\definecolor{success}{RGB}{30,110,60}

\hypersetup{
    colorlinks=true,
    linkcolor=dressmeblue,
    urlcolor=dressmeblue,
    pdftitle={DressMe -- Objectif 4 -- Fiche technique},
    pdfauthor={Projet DressMe}
}

\lstset{
    basicstyle=\ttfamily\small,
    backgroundcolor=\color{lightgray},
    frame=single,
    rulecolor=\color{gray},
    breaklines=true,
    columns=fullflexible,
    showstringspaces=false
}

\pagestyle{fancy}
\fancyhf{}
\lhead{\textit{DressMe -- Objectif 4}}
\rhead{\textit{Fiche technique}}
\cfoot{\thepage}

\titleformat{\section}
    {\Large\bfseries\color{dressmeblue}}
    {\thesection}{0.7em}{}
\titleformat{\subsection}
    {\large\bfseries\color{dressmeblue}}
    {\thesubsection}{0.7em}{}

\newcommand{\pathcode}[1]{\texttt{\path{#1}}}
\newcommand{\done}{\textcolor{success}{\textbf{Réalisé}}}
\newcommand{\partialdone}{\textcolor{warning}{\textbf{Partiel}}}
\newcommand{\notdone}{\textcolor{red}{\textbf{Non réalisé}}}

\begin{document}

\begin{titlepage}
    \centering
    \vspace*{1.5cm}
    {\Huge\bfseries DressMe\par}
    \vspace{0.5cm}
    {\LARGE Fiche technique détaillée\par}
    \vspace{0.4cm}
    {\Large Objectif 4 : détection des lacunes de garde-robe\par}
    \vspace{1.2cm}
    \rule{\textwidth}{1pt}
    \vspace{0.6cm}
    {\large Identification des pièces manquantes permettant de maximiser
    le nombre de tenues réalisables\par}
    \vspace{0.6cm}
    \rule{\textwidth}{1pt}
    \vfill
    \begin{tabular}{ll}
        \textbf{Projet :} & DressMe \\
        \textbf{Objectif :} & Wardrobe gap detection \\
        \textbf{Version analysée :} & Pipeline LLM + FashionCLIP \\
        \textbf{Version des prompts :} & \pathcode{gap-agent-v6} \\
        \textbf{Date :} & 29 octobre 2026
    \end{tabular}
    \vfill
    \textit{Cette fiche décrit l'architecture réellement présente dans le
    dossier du projet, ainsi que les éléments encore incomplets.}
\end{titlepage}

\tableofcontents
\newpage

\section{Objectif et périmètre}

\subsection{Objectif fonctionnel}

L'objectif 4 de DressMe consiste à détecter les pièces absentes d'une
garde-robe et à les classer selon leur capacité à augmenter le nombre de
tenues réalisables.

L'objectif officiel est le suivant :

\begin{quote}
\textit{Wardrobe gap detection -- identifying missing pieces that would
maximize the number of possible combinations, useful for strategic
purchasing.}
\end{quote}

En français, le système doit répondre à la question suivante :

\begin{quote}
\textit{Quelle pièce vestimentaire manque dans la garde-robe et permettrait
de créer le plus grand nombre de nouvelles combinaisons valides ?}
\end{quote}

La mesure principale est le gain marginal :

\[
\mathrm{Gain}(c)
=
N(\mathcal{W}\cup\{c\})-N(\mathcal{W})
\]

où :

\begin{itemize}[leftmargin=*]
    \item $\mathcal{W}$ est la garde-robe actuelle ;
    \item $c$ est une pièce candidate absente ;
    \item $N(\mathcal{W})$ est le nombre de tenues valides réalisables ;
    \item $N(\mathcal{W}\cup\{c\})$ est le nombre de tenues après ajout de la
    pièce candidate.
\end{itemize}

Le classement des lacunes est donc fondé sur le nombre de nouvelles tenues
débloquées, et non uniquement sur la similarité visuelle ou sur une
recommandation subjective.

\subsection{Personas concernées}

\begin{itemize}[leftmargin=*]
    \item \textbf{Amine :} éviter les achats inutiles et privilégier la valeur
    créée par une nouvelle pièce ;
    \item \textbf{Yasmine :} trouver rapidement des tenues adaptées à son
    quotidien ;
    \item \textbf{Sami :} construire une garde-robe cohérente et renforcer une
    signature stylistique.
\end{itemize}

Les personas sont actuellement utilisées comme cadres d'explication. Elles ne
modifient pas automatiquement le classement numérique, sauf si un outil
confirme explicitement qu'une personnalisation a été appliquée.

\subsection{Périmètre actuel}

Le système couvre :

\begin{itemize}[leftmargin=*]
    \item la création d'un inventaire visuel à partir d'images ;
    \item l'extraction d'embeddings avec FashionCLIP ;
    \item la génération de pièces candidates par un LLM ;
    \item l'évaluation de la compatibilité entre les pièces existantes et les
    candidates ;
    \item le comptage déterministe des tenues ;
    \item le calcul de bornes de gain lorsque certaines compatibilités sont
    incertaines ;
    \item l'explication textuelle des résultats par un agent conversationnel ;
    \item un rapport statique pré-calculé utilisé uniquement comme solution de
    repli.
\end{itemize}

\section{État global du projet}

\begin{table}[H]
\centering
\caption{État des principales composantes}
\begin{tabularx}{\textwidth}{>{\raggedright\arraybackslash}p{4.2cm} X c}
\toprule
\textbf{Composante} & \textbf{Description} & \textbf{État} \\
\midrule
Inventaire FashionCLIP & Embeddings normalisés et labels visuels provisoires &
\done \\
Détection des doublons & Exclusion des doublons exacts de l'inventaire &
\done \\
Génération d'archetypes & Génération de pièces manquantes génériques par LLM &
\done \\
Jugement de compatibilité & Classification compatible, incompatible ou incertain &
\partialdone \\
Comptage des tenues & Calcul Python basé sur les templates de l'objectif 4 &
\done \\
Calcul du gain marginal & Calcul déterministe et bornes conservative/optimiste &
\done \\
Classement final LLM & Disponible uniquement après complétion des jugements &
\partialdone \\
Agent conversationnel & Réponses avec outils en lecture seule et garde-fous &
\done \\
Rapport statique de secours & Export du notebook validé et borné &
\done \\
Application frontend/backend & Intégration dans une application complète &
\notdone \\
Images d'une vraie garde-robe utilisateur & Données personnelles réelles &
\notdone \\
Évaluation complète de la qualité & Évaluation finale représentative &
\notdone \\
\bottomrule
\end{tabularx}
\end{table}

\subsection{État de l'exécution actuelle}

Les éléments suivants ont été observés dans le projet :

\begin{itemize}[leftmargin=*]
    \item l'échantillon contient environ 200 images provenant de
    Re-PolyVore ;
    \item les images couvrent 20 catégories avec environ 10 images par
    catégorie ;
    \item plusieurs groupes de doublons exacts ont été détectés ;
    \item les doublons sont exclus avant le calcul de l'inventaire ;
    \item FashionCLIP produit des vecteurs normalisés de dimension 512 ;
    \item cinq archetypes de lacunes ont été générés par le LLM ;
    \item 80 jugements de compatibilité ont été réalisés sur 2\,122 jugements
    nécessaires ;
    \item le calcul complet de compatibilité n'est donc pas terminé ;
    \item aucun classement final complet ne doit être présenté comme définitif
    tant que les 2\,122 jugements ne sont pas disponibles.
\end{itemize}

La progression actuelle est donc :

\[
\frac{80}{2122}\approx 3{,}77\%
\]

Cette valeur décrit la progression technique des jugements, et non une
performance de recommandation.

\section{Architecture générale}

\subsection{Vue d'ensemble}

L'architecture distingue clairement :

\begin{enumerate}[leftmargin=*]
    \item la perception visuelle ;
    \item la génération de candidates ;
    \item le jugement de compatibilité ;
    \item le calcul numérique déterministe ;
    \item l'explication conversationnelle.
\end{enumerate}

Le LLM ne doit pas calculer lui-même les gains. Il fournit des catégories
structurées et des jugements de compatibilité. Le comptage des tenues et le
classement numérique sont effectués par Python.

\begin{figure}[H]
\centering
\begin{tabular}{c}
\colorbox{lightblue}{\parbox{0.75\textwidth}{\centering
Images de la garde-robe utilisateur}} \\[0.2cm]
$\downarrow$ \\[0.2cm]
\colorbox{lightblue}{\parbox{0.75\textwidth}{\centering
FashionCLIP : embeddings et attributs visuels provisoires}} \\[0.2cm]
$\downarrow$ \\[0.2cm]
\colorbox{lightblue}{\parbox{0.75\textwidth}{\centering
LLM : génération d'archetypes de pièces manquantes}} \\[0.2cm]
$\downarrow$ \\[0.2cm]
\colorbox{lightblue}{\parbox{0.75\textwidth}{\centering
LLM : jugements de compatibilité par paire}} \\[0.2cm]
$\downarrow$ \\[0.2cm]
\colorbox{lightblue}{\parbox{0.75\textwidth}{\centering
Python : comptage des tenues et calcul des gains}} \\[0.2cm]
$\downarrow$ \\[0.2cm]
\colorbox{lightblue}{\parbox{0.75\textwidth}{\centering
Agent : classement, justification et réponse utilisateur}}
\end{tabular}
\caption{Flux principal de l'objectif 4}
\end{figure}

\subsection{Modules principaux}

\begin{table}[H]
\centering
\caption{Fichiers impliqués dans l'implémentation}
\begin{tabularx}{\textwidth}{>{\ttfamily\raggedright\arraybackslash}p{5.5cm} X}
\toprule
\textbf{Fichier} & \textbf{Responsabilité} \\
\midrule
\pathcode{dressme\_gap\_agent/fashion\_encoder.py} &
Extraction des embeddings FashionCLIP et découverte des images. \\
\pathcode{dressme\_gap\_agent/compatibility.py} &
Génération des archetypes, préparation des paires et validation des jugements. \\
\pathcode{dressme\_gap\_agent/llm\_gap\_pipeline.py} &
Orchestration de la génération, du cache, des checkpoints et du calcul final. \\
\pathcode{dressme\_gap\_agent/outfit\_engine.py} &
Templates de tenues, comptage des combinaisons et calcul des bornes de gain. \\
\pathcode{dressme\_gap\_agent/sources.py} &
Sources d'inventaire et analyse LLM exposées à l'agent. \\
\pathcode{dressme\_gap\_agent/report.py} &
Chargement et validation du rapport statique de secours. \\
\pathcode{dressme\_gap\_agent/agent.py} &
Agent conversationnel, outils, appels LLM et réponse finale. \\
\pathcode{dressme\_gap\_agent/prompts.py} &
Prompt système, personas, politique de grounding et contexte structuré. \\
\pathcode{dressme\_gap\_agent/config.py} &
Configuration du routeur, des modèles, des limites et des timeouts. \\
\pathcode{dressme\_gap\_agent/evaluation.py} &
Évaluation automatisée des choix d'outils et des réponses. \\
\pathcode{tests/test\_dressme\_gap\_agent.py} &
Tests unitaires et tests de sécurité fonctionnelle. \\
\bottomrule
\end{tabularx}
\end{table}

\section{Prérequis}

\subsection{Environnement logiciel}

Le projet utilise principalement :

\begin{itemize}[leftmargin=*]
    \item Python ;
    \item PyTorch ;
    \item Transformers ;
    \item FashionCLIP ;
    \item NumPy ;
    \item Pillow ;
    \item OpenAI Python SDK ;
    \item python-dotenv ;
    \item unittest pour les tests.
\end{itemize}

Un environnement virtuel est déjà présent dans le projet sous
\pathcode{.venv}.

\subsection{Configuration du modèle}

La configuration est lue depuis des variables d'environnement. Les secrets ne
doivent pas être écrits dans le code source.

\begin{lstlisting}[language=bash]
DRESSME_LLM_BASE_URL=http://127.0.0.1:20128/v1
DRESSME_LLM_API_KEY=<cle-locale>
DRESSME_LLM_MODEL=GapDetection
DRESSME_LLM_COMPAT_MODEL=groq/openai/gpt-oss-120b
DRESSME_LLM_TEMPERATURE=0.2
DRESSME_LLM_MAX_TOKENS=600
DRESSME_LLM_JUDGE_MAX_TOKENS=1600
DRESSME_LLM_TIMEOUT=30
DRESSME_LLM_MAX_RETRIES=1
DRESSME_LLM_BATCH_DELAY=0
\end{lstlisting}

\begin{table}[H]
\centering
\caption{Paramètres du pipeline}
\begin{tabular}{lll}
\toprule
\textbf{Paramètre} & \textbf{Valeur actuelle} & \textbf{Rôle} \\
\midrule
Température & 0,2 & Réduire la variabilité des réponses \\
Tokens réponse & 600 & Limiter l'explication finale \\
Tokens juge & 1600 & Limiter les lots de jugements \\
Timeout & 30 s & Éviter les appels bloqués \\
Retries applicatifs & 1 & Réessayer les erreurs transitoires \\
Taille d'un lot & 10 & Limiter les jugements par appel \\
Max candidats & 6 & Limiter les archetypes générés \\
\bottomrule
\end{tabular}
\end{table}

Ces valeurs sont des configurations initiales. Elles ne constituent pas encore
des hyperparamètres optimisés expérimentalement.

\section{Étape 1 : acquisition et préparation des images}

\subsection{Source actuelle}

Les images actuellement disponibles proviennent d'un échantillon de
Re-PolyVore. Elles ne représentent pas encore une garde-robe personnelle
vérifiée.

L'exécution actuelle vérifie principalement la faisabilité technique du
pipeline. Elle ne permet pas encore de conclure sur la pertinence des lacunes
pour un utilisateur réel.

\subsection{Préparation}

\begin{enumerate}[leftmargin=*]
    \item Découverte des images dans le répertoire de garde-robe.
    \item Association des images à des identifiants.
    \item Vérification de la lisibilité des images.
    \item Encodage par FashionCLIP.
    \item Conversion en vecteurs normalisés.
    \item Filtrage des doublons exacts.
\end{enumerate}

\begin{lstlisting}[language=Python]
load_inventory_items(
    inventory_path,
    exclude_exact_duplicates=True
)
\end{lstlisting}

\subsection{Limites actuelles}

\begin{itemize}[leftmargin=*]
    \item les images personnelles réelles ne sont pas encore intégrées ;
    \item la qualité de la segmentation du vêtement n'est pas évaluée sur un
    jeu de données utilisateur ;
    \item les attributs visuels produits par FashionCLIP sont provisoires ;
    \item une erreur de catégorisation visuelle peut influencer les jugements
    de compatibilité ultérieurs.
\end{itemize}

\section{Étape 2 : extraction des représentations visuelles}

\subsection{FashionCLIP}

FashionCLIP transforme chaque image vestimentaire en un vecteur numérique :

\[
\mathbf{x}_i \in \mathbb{R}^{512}
\]

Le vecteur est normalisé afin de permettre la comparaison par similarité
cosinus :

\[
\mathrm{sim}(\mathbf{x}_i,\mathbf{x}_j)
=
\frac{\mathbf{x}_i\cdot\mathbf{x}_j}
{\|\mathbf{x}_i\|_2\|\mathbf{x}_j\|_2}
\]

FashionCLIP est utilisé pour :

\begin{itemize}[leftmargin=*]
    \item représenter les articles ;
    \item produire des labels visuels provisoires ;
    \item fournir des informations compactes au pipeline de compatibilité ;
    \item éviter d'envoyer systématiquement les images brutes au LLM.
\end{itemize}

FashionCLIP ne calcule pas directement le gain d'une nouvelle pièce.

\subsection{Attributs visuels provisoires}

\begin{itemize}[leftmargin=*]
    \item catégorie ;
    \item couleur ;
    \item motif ;
    \item style ;
    \item slot vestimentaire.
\end{itemize}

Ces attributs servent de contexte au LLM juge. Ils ne constituent pas à eux
seuls une preuve de compatibilité.

\section{Étape 3 : génération des pièces candidates}

\subsection{Principe}

Le système ne cherche pas directement un produit dans un catalogue commercial.
Il demande au LLM de proposer des archetypes génériques de pièces qui pourraient
combler les lacunes de la garde-robe.

Un archetype est décrit par :

\begin{itemize}[leftmargin=*]
    \item un identifiant ;
    \item un slot ;
    \item une catégorie ;
    \item une description ;
    \item une justification stylistique.
\end{itemize}

\begin{lstlisting}[language=json]
{
  "item_id": "candidate-001",
  "slot": "outer",
  "category": "outerwear",
  "description": "veste legere neutre",
  "rationale": "ajoute une couche compatible avec plusieurs tenues"
}
\end{lstlisting}

Les candidates sont des hypothèses de pièces manquantes. Elles ne sont pas des
produits disponibles à l'achat, des liens commerciaux, des recommandations de
prix ou des résultats personnalisés par budget sans prix vérifié.

\section{Étape 4 : préparation des paires de compatibilité}

Les articles existants et les candidates sont regroupés afin de générer les
paires :

\[
(\text{article existant},\text{candidate})
\]

Les paires entre deux candidates ne sont pas utilisées dans le calcul standard
du gain marginal.

Chaque paire possède un identifiant déterministe :

\begin{lstlisting}[language=Python]
pair_id = item_id_left + "::" + item_id_right
\end{lstlisting}

Le pipeline vérifie :

\begin{itemize}[leftmargin=*]
    \item que chaque identifiant attendu est présent ;
    \item qu'aucune paire n'est retournée deux fois ;
    \item qu'aucune paire inconnue n'est acceptée ;
    \item que le verdict appartient à l'ensemble autorisé.
\end{itemize}

\section{Étape 5 : jugement LLM de compatibilité}

\subsection{Verdicts possibles}

\begin{itemize}[leftmargin=*]
    \item \texttt{compatible} ;
    \item \texttt{incompatible} ;
    \item \texttt{uncertain}.
\end{itemize}

Le verdict \texttt{uncertain} est conservé explicitement. Il ne doit pas être
transformé silencieusement en compatibilité ou en incompatibilité.

\subsection{Structure d'un jugement}

\begin{lstlisting}[language=json]
{
  "pair_id": "item-001::candidate-001",
  "verdict": "compatible",
  "evidence": "la couleur neutre et la categorie permettent plusieurs associations"
}
\end{lstlisting}

\subsection{Traitement par lots}

Les jugements sont envoyés par lots de dix paires :

\[
B=10
\]

Après chaque lot, les jugements validés sont enregistrés, les paires déjà
traitées ne sont pas renvoyées et le processus peut être relancé.

\subsection{État incomplet}

L'exécution actuelle contient :

\[
80 \text{ jugements réalisés sur } 2122 \text{ jugements requis}
\]

La génération du classement final est désactivée tant que le nombre de paires
restantes n'est pas nul.

\section{Étape 6 : gestion du cache et reprise}

Les résultats intermédiaires sont conservés dans
\pathcode{wardrobe/.dressme-cache/}.

\begin{itemize}[leftmargin=*]
    \item \pathcode{gap\_archetypes.json} : candidates générées ;
    \item \pathcode{compatibility\_judgments.json} : jugements par paire ;
    \item \pathcode{llm\_gap\_analysis.json} : résultat global du pipeline ;
    \item \pathcode{manifest.json} : informations sur l'inventaire.
\end{itemize}

Le cache est associé à un hash de l'inventaire, une version du prompt, le
modèle configuré et une version du pipeline.

\begin{lstlisting}[language=powershell]
python -B -m dressme_gap_agent.llm_gap_pipeline --inventory wardrobe/.dressme-cache/full/wardrobe_inventory.json --cache-dir wardrobe/.dressme-cache/llm-gap-v1 --output wardrobe/.dressme-cache/llm_gap_analysis.json --batch-size 10 --max-new-pairs 40
\end{lstlisting}

La commande doit être relancée jusqu'à obtenir :

\begin{lstlisting}[language=json]
{
  "status": "complete",
  "ranking_ready": true,
  "compatibility_progress": {
    "remaining_pairs": 0
  }
}
\end{lstlisting}

\section{Étape 7 : définition des templates de tenues}

Le moteur repose sur les templates définis pour l'objectif 4. Les slots
utilisés comprennent :

\begin{itemize}[leftmargin=*]
    \item \texttt{top} ;
    \item \texttt{bottom} ;
    \item \texttt{onepiece} ;
    \item \texttt{outer} ;
    \item \texttt{shoes} ;
    \item \texttt{bag} ;
    \item \texttt{accessory}.
\end{itemize}

\begin{table}[H]
\centering
\caption{Exemples de structures de tenues}
\begin{tabular}{ll}
\toprule
\textbf{Template} & \textbf{Slots possibles} \\
\midrule
Tenue de base & top + bottom + shoes \\
Tenue avec couche & top + bottom + outer + shoes \\
Tenue avec accessoires & top + bottom + shoes + bag + accessory \\
Tenue une pièce & onepiece + shoes \\
Tenue une pièce complète & onepiece + outer + shoes + bag \\
\bottomrule
\end{tabular}
\end{table}

Les templates exacts du notebook restent la source de vérité numérique.

\section{Étape 8 : comptage déterministe des tenues}

Le moteur calcule d'abord :

\[
N(\mathcal{W})
\]

Puis, pour chaque candidate $c$ :

\[
N(\mathcal{W}\cup\{c\})
\]

Le gain est :

\[
\mathrm{Gain}(c)
=
N(\mathcal{W}\cup\{c\})-N(\mathcal{W})
\]

Ce calcul est effectué par Python et non par le LLM. Cette séparation rend les
nombres reproductibles et permet de tester le moteur indépendamment de la
génération de texte.

\section{Étape 9 : gestion de l'incertitude}

Lorsque certaines paires sont incertaines, le moteur calcule :

\[
\mathrm{Gain}_{\mathrm{lower}}(c)
\]

et

\[
\mathrm{Gain}_{\mathrm{upper}}(c)
\]

On obtient :

\[
\mathrm{Gain}_{\mathrm{lower}}(c)
\leq
\mathrm{Gain}_{\mathrm{réel}}(c)
\leq
\mathrm{Gain}_{\mathrm{upper}}(c)
\]

Ces bornes sont des estimations dépendantes des jugements du LLM. Elles ne sont
pas des probabilités calibrées.

\section{Étape 10 : classement des lacunes}

Lorsque tous les jugements sont disponibles, les candidates sont triées par
gain conservateur puis par gain optimiste :

\[
\mathrm{clé}(c)
=
\left(
\mathrm{Gain}_{\mathrm{lower}}(c),
\mathrm{Gain}_{\mathrm{upper}}(c)
\right)
\]

Un résultat peut avoir la forme suivante :

\begin{lstlisting}[language=json]
{
  "item_id": "candidate-001",
  "description": "veste legere neutre",
  "category": "outerwear",
  "slot": "outer",
  "llm_rationale": "ajoute une couche polyvalente",
  "gain_lower": 12,
  "gain_upper": 18,
  "rank_by_conservative_gain": 1
}
\end{lstlisting}

Le classement n'est pas exposé comme prêt tant que :

\begin{lstlisting}[language=json]
"ranking_ready": true
\end{lstlisting}

n'est pas présent.

\section{Étape 11 : rapport statique de secours}

Le fichier \pathcode{gap\_report.json} contient un export pré-calculé du
notebook de l'objectif 4. Il est utilisé uniquement lorsqu'une analyse LLM
complète n'est pas disponible.

\begin{lstlisting}[language=json]
{
  "source": "precomputed_notebook_export",
  "recalculated": false
}
\end{lstlisting}

Le chargeur vérifie le format JSON, la présence des champs obligatoires,
l'unicité des identifiants et la validité des gains. Les chemins locaux
d'images ne sont pas envoyés au LLM.

Le rapport statique ne recalcule pas les recommandations et n'applique pas
automatiquement la persona, la saison ou le budget.

\section{Étape 12 : agent conversationnel}

L'agent explique les résultats au lieu de les recalculer. Il peut résumer
l'inventaire, présenter les lacunes, comparer les candidates et signaler les
contraintes non appliquées.

\begin{table}[H]
\centering
\caption{Outils disponibles}
\begin{tabularx}{\textwidth}{>{\ttfamily\raggedright\arraybackslash}p{4cm} X}
\toprule
\textbf{Outil} & \textbf{Fonction} \\
\midrule
\texttt{get\_wardrobe\_summary} &
Retourner un résumé borné de l'inventaire. \\
\texttt{rank\_wardrobe\_gaps} &
Retourner le classement lorsque l'analyse est complète. \\
\texttt{get\_gap\_report} &
Retourner une tranche du rapport statique pré-calculé. \\
\bottomrule
\end{tabularx}
\end{table}

L'agent ne dispose pas d'outil utilisateur permettant d'exécuter du shell,
d'accéder à une base SQL, de rechercher dans des magasins, d'exécuter du code
arbitraire ou de modifier le classement.

\section{Étape 13 : contexte utilisateur}

Le système accepte un contexte structuré limité :

\begin{lstlisting}[language=json]
{
  "persona": "amine",
  "season": "cold",
  "budget_max_tnd": 250,
  "wardrobe_counts": {
    "top": 3,
    "shoes": 1
  }
}
\end{lstlisting}

Les champs sont validés par liste blanche, validation des types, validation
des bornes numériques et normalisation.

Le contexte est actuellement un contexte d'explication. Il ne modifie pas
automatiquement le classement du moteur. Le budget ne peut pas être vérifié
sans prix réels, et la saison ne modifie pas le rapport statique.

\section{Étape 14 : sécurité et fiabilité}

\begin{itemize}[leftmargin=*]
    \item taille maximale des messages utilisateur ;
    \item validation des arguments d'outils ;
    \item limites sur \texttt{top\_k} et \texttt{limit} ;
    \item validation des chemins de rapport ;
    \item vérification des types JSON ;
    \item rejet des champs inattendus ;
    \item validation des identifiants de paires ;
    \item vérification des verdicts autorisés.
\end{itemize}

Par défaut, le routeur doit être local :
\pathcode{http://127.0.0.1:20128/v1}.

Une route distante est refusée par défaut. Si elle est explicitement autorisée,
elle doit utiliser HTTPS.

Les textes utilisateurs, descriptions de vêtements et résultats d'outils sont
traités comme des données. Le prompt interdit notamment la divulgation des
clés API, chemins locaux, variables d'environnement, prompts internes et
raisonnements internes détaillés.

\section{Étape 15 : évaluation automatisée}

\begin{lstlisting}[language=powershell]
python -B -m dressme_gap_agent.evaluation --report .\gap_report.json --output .\llm_eval_results.json
\end{lstlisting}

Les tests vérifient la sélection des outils, leurs arguments, les identifiants
retournés, les gains retournés, les termes obligatoires et le comportement face
à une tentative d'injection.

La dernière exécution textuelle enregistrée indique :

\begin{itemize}[leftmargin=*]
    \item jeu de données : \texttt{dressme-gap-eval-v2} ;
    \item version du prompt : \texttt{gap-agent-v5} ;
    \item cas automatisés réussis : 6 sur 7 ;
    \item latence moyenne : environ 1,44 seconde ;
    \item coût : non estimé ;
    \item contrôle de non-divulgation face à l'injection : réussi ;
    \item tâche mélangée avec injection : échec partiel.
\end{itemize}

Cette évaluation est antérieure à \texttt{gap-agent-v6} et doit être relancée.

\section{Étape 16 : tests unitaires}

\begin{lstlisting}[language=powershell]
python -m unittest discover -s tests -v
\end{lstlisting}

Les tests couvrent notamment la configuration, les limites des retries, la
construction du prompt, la validation du contexte, le rapport statique, les
paires de compatibilité, le calcul des tenues, les bornes de gain, la
préparation des images et le comportement de l'agent.

\section{Étape 17 : procédure de reproduction complète}

\subsection{Phase A -- vérifier l'environnement}

\begin{lstlisting}[language=powershell]
.\.venv\Scripts\Activate.ps1
python --version
Test-Path .\wardrobe\.dressme-cache\full\wardrobe_inventory.json
\end{lstlisting}

\subsection{Phase B -- exécuter les tests}

\begin{lstlisting}[language=powershell]
python -m unittest discover -s tests -v
\end{lstlisting}

\subsection{Phase C -- préparer l'inventaire}

\begin{enumerate}[leftmargin=*]
    \item Placer les images dans le répertoire de garde-robe.
    \item Vérifier qu'elles sont lisibles.
    \item Exclure les doublons exacts.
    \item Générer les embeddings FashionCLIP.
    \item Vérifier la dimension des embeddings.
    \item Vérifier le nombre de catégories détectées.
\end{enumerate}

\subsection{Phase D -- générer les archetypes}

Vérifier la présence d'au moins une candidate, d'un slot valide, d'une
description et d'une justification, sans lien commercial inventé.

\subsection{Phase E -- exécuter les jugements}

Relancer le pipeline par lots jusqu'à obtenir
\texttt{remaining\_pairs = 0}. Vérifier que les jugements précédents sont
conservés et que les paires déjà traitées ne sont pas renvoyées.

\subsection{Phase F -- calculer le classement}

Lorsque l'analyse est complète, vérifier que
\texttt{ranking\_ready=true}, puis contrôler les gains inférieur et supérieur
de plusieurs candidates.

\subsection{Phase G -- tester l'agent}

Tester les requêtes suivantes :

\begin{enumerate}[leftmargin=*]
    \item « Quelle pièce manque le plus à ma garde-robe ? »
    \item « Compare les deux premières candidates. »
    \item « Quel est le gain de cette veste ? »
    \item « Je suis Amine, quelle recommandation évite le gaspillage ? »
    \item « Quel article puis-je acheter avec 200 TND ? »
    \item une requête contenant une tentative de prompt injection.
\end{enumerate}

\section{Techniques utilisées}

\begin{table}[H]
\centering
\caption{Techniques principales}
\begin{tabularx}{\textwidth}{>{\raggedright\arraybackslash}p{4.2cm} X X}
\toprule
\textbf{Technique} & \textbf{Fonction} & \textbf{Pourquoi} \\
\midrule
FashionCLIP & Représentation vectorielle des images &
Comparer et résumer les articles. \\
Similarité cosinus & Comparaison d'embeddings &
Mesurer la proximité visuelle. \\
Zero-shot labels & Labels visuels sans classification supervisée locale &
Obtenir des catégories provisoires. \\
LLM structured output & Archetypes et verdicts JSON &
Imposer une structure vérifiable. \\
LLM judging & Jugement de compatibilité selon une consigne &
Rendre les critères plus homogènes. \\
Caching & Sauvegarde des résultats intermédiaires &
Éviter de refaire les appels réussis. \\
Checkpointing & Reprise après interruption &
Permettre une exécution par lots. \\
Comptage combinatoire & Énumération des tenues valides &
Calculer le gain de façon déterministe. \\
Bornes d'incertitude & Gains inférieur et supérieur &
Ne pas masquer les jugements inconnus. \\
Tool calling & Accès contrôlé aux données &
Séparer les calculs de l'explication. \\
Validation JSON & Contrôle des sorties du modèle &
Réduire les résultats malformés. \\
\bottomrule
\end{tabularx}
\end{table}

\section{Limites actuelles}

\subsection{Limites des données}

\begin{itemize}[leftmargin=*]
    \item les images sont issues d'un échantillon externe ;
    \item la garde-robe personnelle d'un utilisateur n'a pas encore été validée ;
    \item les attributs visuels ne sont pas encore évalués avec une vérité
    terrain complète.
\end{itemize}

\subsection{Limites du modèle de compatibilité}

\begin{itemize}[leftmargin=*]
    \item les verdicts de compatibilité sont produits par un LLM ;
    \item les jugements ne sont pas calibrés sur suffisamment d'annotations
    humaines ;
    \item le résultat dépend du modèle et du prompt ;
    \item la stabilité entre plusieurs fournisseurs n'est pas démontrée ;
    \item l'analyse actuelle est incomplète.
\end{itemize}

\subsection{Limites du classement}

\begin{itemize}[leftmargin=*]
    \item aucun classement complet ne doit être annoncé avant la fin des
    2\,122 jugements ;
    \item les bornes de gain sont des estimations ;
    \item les gains ne prennent pas encore en compte des prix réels ;
    \item le budget n'est pas intégré au calcul ;
    \item la saison et la météo ne sont pas intégrées au moteur ;
    \item une candidate est un archetype et non un produit achetable.
\end{itemize}

\subsection{Limites de l'application}

\begin{itemize}[leftmargin=*]
    \item aucune interface frontend complète n'est présente dans ce dossier ;
    \item aucune intégration backend de production n'est finalisée ;
    \item la gestion des comptes utilisateurs n'est pas implémentée ;
    \item les préférences persistantes ne sont pas encore connectées ;
    \item la confidentialité et la rétention du fournisseur externe restent à
    vérifier.
\end{itemize}

\section{Travail restant}

\begin{longtable}{p{4.8cm} p{7.2cm} p{2.3cm}}
\caption{Éléments à compléter avant une version finale} \\
\toprule
\textbf{Tâche} & \textbf{Description} & \textbf{Priorité} \\
\midrule
\endfirsthead
\toprule
\textbf{Tâche} & \textbf{Description} & \textbf{Priorité} \\
\midrule
\endhead
Terminer les jugements & Exécuter les 2\,042 paires restantes et obtenir
\texttt{remaining\_pairs = 0}. & Haute \\
Valider le classement & Vérifier manuellement les candidates et les gains
calculés. & Haute \\
Utiliser des images personnelles & Tester l'inventaire sur une vraie
garde-robe utilisateur. & Haute \\
Créer une vérité terrain & Faire annoter les compatibilités par des personnes
ou des experts. & Haute \\
Évaluer la qualité du gap detection & Mesurer précision, rappel, stabilité et
corrélation avec les avis humains. & Haute \\
Mettre à jour l'évaluation LLM & Rejouer le jeu d'évaluation avec
\texttt{gap-agent-v6}. & Haute \\
Ajouter les prix & Associer les archetypes à un catalogue fiable et local. &
Moyenne \\
Intégrer le budget & Recalculer le classement sous contrainte de prix. &
Moyenne \\
Intégrer la saison & Ajouter météo, saison et occasion dans les templates. &
Moyenne \\
Tester plusieurs modèles & Comparer qualité, coût et stabilité. & Moyenne \\
Tester la vision distante & Vérifier l'envoi d'images vers le fournisseur
configuré. & Moyenne \\
Construire l'interface & Relier le pipeline à une application utilisable par
l'utilisateur final. & Moyenne \\
Évaluer la confidentialité & Documenter la conservation et le traitement des
données par le fournisseur. & Haute \\
Étudier le fine-tuning & Préparer un jeu de données annoté et un modèle dont
les poids sont contrôlables. & Faible \\
\bottomrule
\end{longtable}

\section{Décision concernant le fine-tuning}

Aucun fine-tuning n'a été réalisé dans l'état actuel du projet.

Cette décision est cohérente avec l'état de l'objectif 4 :

\begin{itemize}[leftmargin=*]
    \item le problème principal est actuellement la complétude et la
    validation des jugements ;
    \item le calcul numérique est déjà séparé du LLM ;
    \item les prompts et les outils sont encore en phase d'évaluation ;
    \item les données annotées sont insuffisantes pour entraîner un modèle ;
    \item le modèle hébergé par le routeur ne fournit pas nécessairement ses
    poids.
\end{itemize}

\section{Synthèse technique}

L'objectif 4 de DressMe repose sur une architecture hybride :

\begin{enumerate}[leftmargin=*]
    \item FashionCLIP encode les articles de la garde-robe ;
    \item un LLM propose des archetypes de pièces manquantes ;
    \item un LLM juge les compatibilités entre articles ;
    \item Python applique les règles de composition des tenues ;
    \item Python calcule les gains et leurs bornes ;
    \item un agent explique les résultats avec des outils contrôlés.
\end{enumerate}

\begin{quote}
\textbf{Le LLM propose et explique ; le code valide, compte et classe.}
\end{quote}

L'implémentation actuelle est fonctionnelle sur le plan architectural, mais
elle ne constitue pas encore une validation finale de l'objectif, car :

\begin{itemize}[leftmargin=*]
    \item l'analyse de compatibilité n'est pas terminée ;
    \item les images personnelles ne sont pas encore utilisées ;
    \item le classement final complet n'est pas encore disponible ;
    \item l'évaluation humaine n'est pas encore réalisée ;
    \item le budget, la saison et les prix ne sont pas intégrés ;
    \item l'application utilisateur complète reste à construire.
\end{itemize}

La prochaine étape prioritaire est de terminer les jugements de compatibilité,
de produire un fichier d'analyse avec \texttt{ranking\_ready=true}, puis de
valider les résultats sur une véritable garde-robe et avec des annotations
humaines.

\end{document}
